% SHARED SECTION, \input by both NonGravForces_Theory_and_Implementation.tex and
% NonGravForces_UserGuide.tex. It carries no preamble and opens no document, so that the
% theory exists ONCE and both documents render the same words (R1).
\small
\begin{itemize}[nosep,leftmargin=*]
\item Bandikova T., McCullough C., Kruizinga G.L., Save H., Christophe B. (2019). GRACE accelerometer data transplant. Adv.\ Space Res.\ 64:623--644.
\item Behzadpour S., Mayer-Gürr T., Krauss S. (2021). GRACE Follow-On accelerometer data recovery. JGR Solid Earth 126:e2020JB021297.
\item Bettadpur S. (2012). GRACE product specification document, GRACE 327-720 (CSR-GR-03-02).
\item Bhattarai S. et al.\ (2019). Anisotropic Earth radiation pressure models. J.\ Geodesy.
\item Bowman B.R. et al.\ (2008). JB2008. AIAA 2008-6438.
\item Bruinsma S., Boniface C. (2021). DTM-2020. J.\ Space Weather Space Clim.\ 11:47.
\item Doornbos E. (2011). Thermospheric density and wind determination from satellite dynamics. PhD, TU Delft; Springer 2012.
\item Drob D.P. et al.\ (2015). HWM14. Earth and Space Science 2:301--319.
\item Emmert J.T. et al.\ (2021). NRLMSIS 2.0: a whole-atmosphere empirical model of temperature and neutral species densities. Earth and Space Science 8:e2020EA001321. \emph{The toolbox ships NRLMSIS 2.1, whose source is the same distribution; this is the model paper.}
\item Flury J., Bettadpur S., Tapley B.D. (2008). Precise accelerometry onboard GRACE. Adv.\ Space Res.\ 42:1414--1423.
\item Guo X., Zhang R. (2024). Performance of NRLMSIS 2.1 using GRACE-A and Swarm-C. Adv.\ Space Res.\ 74:2475.
\item Hu X. et al.\ (2022). Estimation of the centre of mass of GRACE-type satellites. Remote Sensing 14:4030.
\item IAGA (2024). IGRF-14. Zenodo 10.5281/zenodo.14012303.
\item Killick R., Fearnhead P., Eckley I.A. (2012). Optimal detection of changepoints with a linear computational cost. JASA 107:1590--1598.
\item Klinger B., Mayer-Gürr T. (2016). Accelerometer calibration within GRACE gravity field recovery. Adv.\ Space Res.\ 58:1597--1609.
\item Knocke P.C., Ries J.C., Tapley B.D. (1988). Earth radiation pressure effects on satellites. AIAA 88-4292.
\item Kornfeld R.P. et al.\ (2019). GRACE-FO: the extension of the legacy. J.\ Spacecraft Rockets.
\item Lévèque T. et al.\ (2021, 2022). Quantum accelerometers for space gravimetry; CARIOQA definition; rotation effects. J.\ Geodesy 95:15; arXiv:2211.01215; arXiv:2212.06686.
\item March G., Doornbos E.N., Visser P.N.A.M. (2019). High-fidelity geometry models for improving the consistency of CHAMP, GRACE, GOCE and Swarm thermospheric density data sets. Adv.\ Space Res.\ 63:213--238. DOI 10.1016/j.asr.2018.07.009. \emph{Geometry, not gas--surface interaction.}
\item March G., Visser P.N.A.M., Doornbos E.N., Siemes C. (2019). CHAMP and GOCE thermospheric wind characterization with improved gas-surface interactions modelling. Adv.\ Space Res.\ 64:1225--1242. DOI 10.1016/j.asr.2019.06.023.
\item March G., van den IJssel J., Siemes C., Visser P.N.A.M., Doornbos E.N., Pilinski M. (2021). Gas-surface interactions modelling influence on satellite aerodynamics and thermosphere mass density. J.\ Space Weather Space Clim.\ 11:54. DOI 10.1051/swsc/2021035. \emph{The fitted $\alpha_E=0.85$ every preset uses.}
\item McCullough C. et al.\ (2019). Calibrated GRACE-FO accelerometer data products (ACT). JPL D-103863.
\item Montenbruck O., Gill E. (2000). Satellite Orbits. Springer.
\item Moe K., Moe M.M. (2005). Gas--surface interactions and satellite drag coefficients. Planet.\ Space Sci.\ 53:793.
\item Peterseim N. et al.\ (2012). Magnetic torquer induced disturbing signals within GRACE accelerometer data. Adv.\ Space Res.
\item Picone J.M. et al.\ (2002). NRLMSISE-00. JGR Space Physics 107:1468.
\item Pilinski M.D., Argrow B.M., Palo S.E. (2010). Semi-empirical model for satellite energy-accommodation coefficients. J.\ Spacecraft Rockets 47(6):951--956. DOI 10.2514/1.49330.
\item Pilinski M.D., Argrow B.M., Palo S.E., Bowman B.R. (2013). Semi-empirical satellite accommodation model for spherical and randomly tumbling objects. J.\ Spacecraft Rockets 50(3):556--571. DOI 10.2514/1.A32348. \emph{SESAM; the source of \texttt{drag.accommodation = "sesam"}.}
\item Robertson R.V. (2015). Highly physical solar radiation pressure modelling with atmospheric effects (SOLAARS). PhD, Virginia Tech; Zhang et al.\ (2019).
\item Sentman L.H. (1961). Free molecule flow theory and its application to the determination of aerodynamic forces. Lockheed LMSC-448514.
\item Siemes C. et al.\ (2023). New thermosphere neutral mass density and crosswind datasets from CHAMP, GRACE and GRACE-FO. J.\ Space Weather Space Clim.\ 13:16.
\item van Helleputte T. (2009). GPS-based orbit determination and accelerometer calibration for LEO satellites. PhD, TU Delft.
\item Vielberg K., Kusche J. (2020). Extended forward and inverse modeling of radiation pressure accelerations for LEO satellites. J.\ Geodesy 94:43. DOI 10.1007/s00190-020-01368-6.
\item Wen H.Y. et al.\ (2019). GRACE-FO Level-1 data product user handbook. JPL D-56935.
\item Wöske F. (2019/2021). PhD thesis, University of Bremen (ZARM); Wöske F., Kato T., Rievers B., List M. (2019). GRACE accelerometer calibration by high-precision non-gravitational force modelling. Adv.\ Space Res.\ 63:1318--1335.
\item Walker A., Mehta P., Koller J. (2014). Drag coefficient model using the Cercignani--Lampis--Lord gas-surface interaction model. J.\ Spacecraft Rockets 51(5):1544--1563. DOI 10.2514/1.A32677. \emph{The adsorbate coverage law of \texttt{drag.accommodation = "walker2014"}; its closed forms are whole-body and normal-incidence only, which is why \texttt{drag.gsi = "cll"} is refused.}
\item Paul S.N., Sheridan P.L., Licata R.J., Mehta P.M. (2022). Comparison of local and global gas-surface interaction models. arXiv:2210.08364v2. \emph{Restates the Walker coverage law and the Mehta substrate coefficient in the form \texttt{kernel.accommodation} implements.}
\item H\l{}adczuk N.A., van den IJssel J., Kodikara T., Siemes C., Visser P. (2024). GRACE-FO radiation pressure modelling for accurate density and crosswind retrieval. Adv.\ Space Res.\ 73:2355--2373. \emph{Table 3 grades the thermal and geometry choices; the source of the 4.22 and 3.04 nm/s$^2$ figures quoted in the limitations.}
\item Zahzam N. et al.\ (2022). Hybrid electrostatic--atomic accelerometer for future space gravity missions. Remote Sensing 14:3273.
\end{itemize}
